\chapopen{1}

\parstart{1.1}{むかしむかし、あるところ}
\parnum{1.1}
\begin{frank}
\chr{}{むかしむかし、}{むかしむかし}{mukashi mukashi,}{\dw{昔}{mukashi} long ago, the old days; said twice to open a tale}{once upon a time,}
\chr{}{あるところに、}{あるところに}{aru tokoro ni,}{\dw{ある}{aru} a certain; \dw{所}{tokoro} place; \pw{に} in, at}{in a certain place,}
\chr{}{おじいさんと}{おじいさんと}{ojiisan to}{\dw{お爺さん}{ojiisan} old man, grandfather; \pw{と} and}{an old man and}
\chr{}{おばあさんが}{おばあさんが}{obāsan ga}{\dw{お婆さん}{obāsan} old woman, grandmother; \pw{が} marks the subject}{an old woman}
\chr{}{住んでいました。}{すんでいました}{sunde imashita.}{\vb{住む}{sumu}{住み}{sumi}{住んで}{sunde}{to live (godan; intr.)}; \textit{-te imashita} was doing, polite}{were living.}
\end{frank}

\parnum{1.2}
\begin{frank}
\chr{}{おじいさんは}{おじいさんは}{ojiisan wa}{\pw{は} marks the topic, read \textit{wa}}{the old man}
\chr{}{山へ柴刈りに、}{やまへしばかりに}{yama e shibakari ni,}{\dw{山}{yama} mountain; \pw{へ} to, read \textit{e}; \dw{柴刈り}{shibakari} gathering firewood, \nobreak+\,\pw{に} (going) in order to}{to the mountain to cut firewood,}
\chr{}{おばあさんは}{おばあさんは}{obāsan wa}{}{the old woman}
\chr{}{川へ洗濯に}{かわへせんたくに}{kawa e sentaku ni}{\dw{川}{kawa} river; \dw{洗濯}{sentaku} washing, laundry}{to the river to do the washing}
\chr{}{行きました。}{いきました}{ikimashita.}{\vb{行く}{iku}{行き}{iki}{行って}{itte}{to go (godan; intr.)}; \textit{-mashita} polite past}{went.}
\end{frank}

\parnum{1.3}
\begin{frank}
\chr{}{おばあさんが}{おばあさんが}{obāsan ga}{}{the old woman}
\chr{}{川で洗濯を}{かわでせんたくを}{kawa de sentaku o}{\pw{で} at, the place of an action; \pw{を} marks the object, read \textit{o}}{the washing at the river}
\chr{}{していると、}{していると}{shite iru to,}{\vb{する}{suru}{し}{shi}{して}{shite}{to do (irregular)}; \pw{と} when, as soon as}{as she was doing,}
\chr{}{大きな桃が}{おおきなももが}{ōkina momo ga}{\dw{大きな}{ōkina} big; \dw{桃}{momo} peach}{a big peach}
\chr{}{流れてきました。}{ながれてきました}{nagarete kimashita.}{\vb{流れる}{nagareru}{流れ}{nagare}{流れて}{nagarete}{to flow (ichidan; intr.)}; \textit{-te kuru} to come (doing)}{came floating along.}
\end{frank}

\parend

\parstart{1.2}{おばあさんは桃を拾って、}
\parnum{2.1}
\begin{frank}
\chr{}{おばあさんは}{おばあさんは}{obāsan wa}{}{the old woman}
\chr{}{桃を拾って、}{ももをひろって}{momo o hirotte,}{\vb{拾う}{hirou}{拾い}{hiroi}{拾って}{hirotte}{to pick up (godan; tr.)}}{picked up the peach and}
\chr{}{家に}{いえに}{ie ni}{\dw{家}{ie} house, home}{home}
\chr{}{持って帰りました。}{もってかえりました}{motte kaerimashita.}{\vb{持つ}{motsu}{持ち}{mochi}{持って}{motte}{to hold, to carry (godan; tr.)}; \vb{帰る}{kaeru}{帰り}{kaeri}{帰って}{kaette}{to return (godan; intr.)}}{carried it back.}
\end{frank}

\parnum{2.2}
\begin{frank}
\chr{}{おじいさんと二人で}{おじいさんとふたりで}{ojiisan to futari de}{\dw{二人}{futari} two people; \pw{で} as, being}{together with the old man}
\chr{}{桃を切ってみると、}{ももをきってみると}{momo o kitte miru to,}{\vb{切る}{kiru}{切り}{kiri}{切って}{kitte}{to cut (godan; tr.)}; \textit{-te miru} to try doing}{when they cut the peach open,}
\chr{}{中から}{なかから}{naka kara}{\dw{中}{naka} inside; \pw{から} from}{from inside}
\chr{}{元気な男の子が}{げんきなおとこのこが}{genki na otoko no ko ga}{\dw{元気}{genki} healthy, lively, \nobreak+\,\pw{な} the adjective link; \dw{男の子}{otoko no ko} boy}{a lively boy}
\chr{}{生まれました。}{うまれました}{umaremashita.}{\vb{生まれる}{umareru}{生まれ}{umare}{生まれて}{umarete}{to be born (ichidan; intr.)}}{was born.}
\end{frank}

\parnum{2.3}
\begin{frank}
\chr{}{二人は}{ふたりは}{futari wa}{}{the two of them}
\chr{}{その子を}{そのこを}{sono ko o}{\dw{その}{sono} that; \dw{子}{ko} child}{that child}
\chr{}{桃太郎と名づけて、}{ももたろうとなづけて}{Momotarō to nazukete,}{\dw{桃太郎}{Momotarō} Peach Boy; \vb{名づける}{nazukeru}{名づけ}{nazuke}{名づけて}{nazukete}{to name (ichidan; tr.)}}{named Momotarō, and}
\chr{}{大切に育てました。}{たいせつにそだてました}{taisetsu ni sodatemashita.}{\dw{大切}{taisetsu} precious, with care; \vb{育てる}{sodateru}{育て}{sodate}{育てて}{sodatete}{to raise, to bring up (ichidan; tr.)}}{raised him with great care.}
\end{frank}

\chapend
